% !TEX TS-program = xelatex
% !TEX encoding = UTF-8 Unicode

% 兰州大学研究生学位论文模板（LZUThesis-PhD）主文件
% 本文件只负责“组装”，各部分的实际内容见 data/ 目录：
%   data/setup.tex          封皮信息（学校代码/题目/作者/导师等）
%   data/statement.tex      诚信说明页 / 授权说明书
%   data/abstract.tex       中英文摘要
%   data/chap01.tex         正文第一章（示例），后续章节按 chapNN.tex 继续添加
%   data/achievements.tex   在学期间的研究成果
%   data/acknowledgements.tex 致谢
% 编译方式：xelatex --> biber --> xelatex --> xelatex

\documentclass[AutoFakeBold]{LZUThesis-PhD}

\begin{document}

%===== 封皮（学校代码/题目/作者/导师等）=====%
%======================================================%
% 封皮页填写内容（学校代码/题目/作者/导师等）
%======================================================%
% 推荐用键值方式填写（等价于旧的 \schoolcode{...} 等命令，二者兼容、可混用）。
% 注意：
%   - 中文/英文标题必须用双层花括号，如 title = {{第一行}{第二行}}，
%     外层是 \lzusetup 的值、内层是标题的换行分组；单层括号将导致封面解析失败。
%   - degree-type 博士/硕士；degree 学术学位/专业学位。
\lzusetup{
  school-code = {10730},
  secret      = {公开},
  cid         = {请修改成你自己专业相关的分类号，如23333},
  degree-type = {博士},   % 或 硕士
  degree      = {\quad 学\quad 术\quad 学\quad 位\quad},  % 或 专业学位
  title       = {{中文}{标题}},                            % 标题样式见下方旧命令注释
  title-en    = {{ English}{ Title}},                     % 英文多行开头需加空格防粘连
  author      = {作者姓名},
  major       = {一级学科·专业},                          % 或 专业学位类型·培养领域
  research    = {研究方向},
  education   = {学历教育/同等学力人员申请博士学位},      % 或 学历教育/同等学力人员申请硕士学位
  supervisor  = {xxx 教授/研究员},
  co-supervisor = {xxx 教授/研究员},                       % 合作导师，可为空，但不可没有这一栏
  work-period = {2020 年 7 月\quad 至 \quad 2021 年 3 月},
  defense-date = {2021 年 5 月},
}
% 若想沿用旧命令，等价写法（保留注释供参考）：
% \schoolcode{10730}
% \secret{公开}
% \cid{请修改成你自己专业相关的分类号，如23333}
% \yjsType{博士}
% \yjsZsZy{\quad 学\quad 术\quad 学\quad 位\quad}
% \title{{中文}{标题}}
% \entitle{{ English}{ Title}}
% \author{作者姓名}
% \major{一级学科·专业}
% \research{研究方向}
% \education{学历教育/同等学力人员申请博士学位}
% \advisor{xxx 教授/研究员}
% \codvisor{xxx 教授/研究员}
% \elapse{2020 年 7 月\quad 至 \quad 2021 年 3 月}
% \defense{2021 年 5 月}


\maketitle

%===== 诚信说明页 / 授权说明书 =====%
%======================================================%
% 诚信说明页 / 授权说明书
%======================================================%
% 如果超出边界，可以调整签字的宽度，现在是50，如果你不用，把下面的注释就好

% 你的签名
\mysignature{
    % \raisebox{-5pt}{
    % \includegraphics[width=40pt]{signature.pdf}
    % }
}
% 你手写的日期
\mytime{
    % \raisebox{-5pt}{
    % \includegraphics[width=40pt]{signature.pdf}
    % }
}
% 老师的手写签名
\supervisorsignature{
    % \raisebox{-5pt}{
    % \includegraphics[width=40pt]{signature.pdf}
    % }
}
% 老师手写的时间
\teachertime{
    % \raisebox{-5pt}{
    % \includegraphics[width=40pt]{signature.pdf}
    % }
}
% 老师手写的成绩
\recommendedgrade{
    % \raisebox{-5pt}{
    % \includegraphics[width=40pt]{signature.pdf}
    % }
}

\makestatement


\frontmatter

%===== 中英文摘要 =====%
%======================================================%
% 中英文摘要
%======================================================%

% 中文摘要
\ZhAbstract{空白页 空白页 空白页 空白页 空白页 空白页 空白页 空白页 空白页 空白页 空白页 空白页 空白页 空白页 空白页 空白页 空白页 空白页 空白页 空白页 空白页 空白页 空白页 空白页 空白页 空白页 空白页 空白页 空白页 空白页 空白页 空白页 空白页 空白页 空白页 空白页 空白页 空白页 空白页 空白页 空白页 空白页 空白页 空白页 空白页 空白页 空白页 空白页 空白页 空白页 空白页 空白页 空白页 空白页 空白页 空白页 空白页 空白页 空白页 空白页 空白页 空白页 空白页 空白页 空白页 空白页 空白页 空白页 空白页 空白页 空白页 空白页 空白页 空白页 空白页 空白页 空白页 空白页 空白页 空白页 空白页 空白页 空白页 空白页 空白页}{关键词1，关键词2}

% 英文摘要
\EnAbstract{example example example example example example example example example example example example example example example example example example example example example example example example example example example example example example example example example example example example example example example example example example example example example example example example example example example example example example example example example example example example example example example example example example example example example example example example example example example example example example example example example example example example example .
    % \fontspec{Times New Roman} {Times New Roman}
}{key-word-1,key-word-2}


%===== 目录（含图表目录）=====%
% \tableofcontents
% 下面这个包含图表目录
\customcontent

% % 部分同学需要专业术语注释表，* 表示不加入目录
% \chapter*{专业术语注释表}
% \begin{longtable}{lll}
%   \caption*{缩略词说明}\\
%   SS & Spread Spectrum & 扩展频谱 \\
%   PAPR & Peak to Average Power Ratio & 峰均比\\
%   DCSK & Differential Chaos Shift Keying &差分混移位键控\\
%   dasd & fdhfudw eqwrqw fasfasfs fewev wqfwefew &\tabincell{l}{太长了\\换行一下}\\
% \end{longtable}

%===== 正文 =====%
\mainmatter

\input{data/chap01}

%===== 参考文献 =====%
\backmatter

% 引入参考文献文件
\printbib
% \nocite{*} %显示数据库中有的，但是正文没有引用的文献

%===== 在学期间研究成果 / 致谢 =====%
%======================================================%
% 在学期间的研究成果
%======================================================%

\Achievements
一、发表论文

1.Article here sd

\blank

二、参与课题

%======================================================%
% 致谢
%======================================================%

\Thanks

这里是致谢页。


\end{document}
